\begin{Definition}
\label{def:banach-valued-analytic-function}
设 $G\subset\mathbb{C}$ 为区域, $Y$ 为 Banach 空间. 对于映射
\begin{equation}\label{eq:banach-valued-complex-function-map}
\mathbb{C}\supset G\ni\lambda\mapsto B(\lambda)\in Y
\end{equation}
若以下极限存在, 则称它在点 $\lambda_0\in\mathbb{C}$ 处\textbf{解析}:
\begin{equation}\label{eq:banach-valued-analytic-function-difference-quotient}
\lim_{\lambda\to\lambda_0}\dfrac{B(\lambda)-B(\lambda_0)}{\lambda-\lambda_0}
\end{equation}

\end{Definition}

正如前面所暗示的, 主要结果如下.

\begin{Theorem}
\label{thm:bounded-operator-analytic-resolvent}
设 $A\in\mathcal{L}(X)$. 那么预解算子 $R_\lambda(A)$ (看作 $\lambda$ 的函数) 在预解集 $\rho(A)$ 上解析.
\end{Theorem}

\begingroup
\postdisplaypenalty=10000
\renewcommand{\endProof}{\par\nobreak\hfill$\blacksquare$\par\vspace{0.1cm}}
\begin{Proof}
差商极限的存在性直接来自 Theorem~\ref{thm:local-resolvent-power-series} 中的幂级数表示.
\end{Proof}
\endgroup

现在断言, 标准复变函数课程中对解析函数建立的方法和结果, 可以放心地用于取值于 Banach 空间的解析函数. 对这一思想更深入的阐述, 例如可参见 \cite{DS}. 作为将已有方法用于这一新背景的例子, 下面证明以下结果.

\begin{Theorem}
\label{thm:bounded-operator-nonempty-spectrum}
非零 Banach 空间上的有界算子的谱至少有一个元素.
\end{Theorem}

\begingroup
\postdisplaypenalty=10000
\renewcommand{\endProof}{\par\nobreak\hfill$\blacksquare$\par\vspace{0.1cm}}
\begin{Proof}
设 $A\in\mathcal{L}(X)$, 并假设 $\sigma(A)$ 为空集; 即预解集为整个复平面. 根据 Theorem~\ref{thm:bounded-operator-analytic-resolvent}, 预解算子 $R_\lambda(A)$ (看作 $\lambda$ 的函数) 是整的; 即在整个复平面上解析. 现在注意到, $\lambda\mapsto R_\lambda(A)$ 在整个 $\mathbb{C}$ 上有界. 为看出这一点, 注意由 \eqref{eq:resolvent-series-outside-norm-disk} 可得
\begin{equation}\label{eq:resolvent-large-parameter-bound}
\|R_\lambda(A)\|\leq\dfrac{1}{\|A\|}\quad\text{for }|\lambda|\geq2\|A\|.
\end{equation}
此外, 由于 $\lambda\mapsto R_\lambda(A)$ 解析, 它在任意有界圆盘上必然有界. 因而可利用 Liouville 定理推出 $\lambda\mapsto R_\lambda(A)$ 为常量. 这产生矛盾, 证明完毕.
\end{Proof}
\endgroup

\begin{Remark}
\label{rem:unbounded-operator-empty-spectrum-example}
Theorems~\ref{thm:local-resolvent-power-series} 和 \ref{thm:bounded-operator-analytic-resolvent} 可推广到闭算子 (证明类似, 参见 Problem~\ref{ex:closed-unbounded-resolvent-analyticity}). 然而, 无界算子可能具有空谱. 例如, 设 $X=L^2(0,1)$, 并令
\begin{equation}\label{eq:empty-spectrum-derivative-domain}
\mathcal{D}(S):=\{y\in H^1(0,1)\mid y(0)=0\}
\end{equation}
及
\begin{equation}\label{eq:empty-spectrum-derivative-action}
Sy=i\dfrac{\mathrm{d}y}{\mathrm{d}x}.
\end{equation}
请读者验证, 对任意 $\lambda\in\mathbb{C}$, 由
\begin{equation}\label{eq:empty-spectrum-derivative-explicit-resolvent}
L_\lambda(y)(x):=-i\int_0^xe^{-i\lambda(x-s)}y(s)\,\mathrm{d}s
\end{equation}
给出且定义域为
\begin{equation}\label{eq:empty-spectrum-derivative-resolvent-domain}
\mathcal{D}(L_\lambda):=L^2(0,1)
\end{equation}
的算子 $L_\lambda$ 确实为预解算子 $R_\lambda(S)$.
\end{Remark}
